%%%%%%%%%%%%%%%%%%%%%%%%%%%%%%%%%%%%%%%%%%%%%%%%%%%%%%%%%%%%%%%%%%%%%%%%%%%%%%%%
% AGENT WARNING (2026-09-18) --- do not use this section when drafting earlier
% sections; this block does not enter the PDF.
%
% The whole section, especially the final outlook sentence beginning
% "Future work will further consider...", is a temporary AI-generated
% placeholder. It is not a settled research boundary or a future-work list.
% Do not cite it when writing the abstract, introduction, methods, or
% simulation, and do not infer gaps, contributions, or section organization
% from this outlook.
% The figures 96.7\%/96.0\% are only temporary placeholders for the noisy
% identification results in Section 6.3. Real results are of a similar
% magnitude. The noise-free Section 6.2 figures 94.7\%/94.3\% and related
% results are real experimental results. This warning only asks that the
% noisy-identification numbers and the placeholder outlook be ignored. It
% must not be used to judge that the experiments as a whole have not
% succeeded, that the real results differ greatly, or to weaken the stated
% contributions while writing. The final conclusion should be rewritten
% after the Section 6.3 numbers are calibrated. See the writing AGENT
% WARNING in main.tex.
%%%%%%%%%%%%%%%%%%%%%%%%%%%%%%%%%%%%%%%%%%%%%%%%%%%%%%%%%%%%%%%%%%%%%%%%%%%%%%%%
\section{Conclusion}
This paper addresses the requirement that an evader simultaneously avoid two NCVs and retain terminal accuracy at a specified target, and proposes a three-dimensional integrated evasion and guidance framework. A four-body relative-motion model is established to describe the coupled interaction among the evader, the target, and the two NCVs. On this basis, the dual-NCV evasion requirement, the target-arrival accuracy, and the control-effort penalty are placed in a unified performance index, and an optimal closed-loop guidance law is derived. The guidance law explicitly accounts for cross-coupling: one maneuver simultaneously affects both evasive encounters and terminal arrival at the target.

To select guidance weights adaptively under different engagement geometries, a BC-TD3 parameter-optimization strategy is proposed. Expert demonstrations generated by single shooting are used to pre-train the actor and pre-fill the replay buffer. During TD3 fine-tuning, a gradually decayed behavior-cloning regularization term retains the expert constraint in the early stage and then transfers the actor toward reward-driven optimization, thereby improving its adaptability in a nonlinear multi-constraint environment. In addition, a recursive least-squares identifier with sparse regularization and a forgetting factor estimates the nontrivial parameters of the state-space model from noisy measurements, which reduces the sensitivity of the guidance loop to measurement errors.

Simulation results show that the proposed method can reliably complete evasion and target arrival at the same time in a three-dimensional dual-NCV scenario. In noise-free Monte Carlo trials, the evasion success rate and the target-arrival success rate of BC-TD3 are 94.7\% and 94.3\%, respectively, which are higher than those of classical proportional navigation, programmed sinusoidal maneuvering, and TD3 without behavior cloning. Under measurement noise, introducing sparse online identification further raises the two success rates to 96.7\% and 96.0\%. These results indicate that the proposed framework can provide an effective and robust solution for integrated multi-NCV evasion and terminal guidance. Future work will further consider vehicle dynamics closer to practice, heterogeneous NCVs, actuator constraints, and verification under a wider range of uncertainties.
